\documentclass[aspectratio=169,11pt]{beamer}
\usepackage[T1]{fontenc}
\usepackage{lmodern,amsmath,booktabs,graphicx}
\definecolor{ink}{HTML}{183248}
\definecolor{teal}{HTML}{187B84}
\definecolor{rust}{HTML}{CB593C}
\setbeamercolor{normal text}{fg=ink,bg=white}
\setbeamercolor{frametitle}{fg=ink}
\setbeamercolor{structure}{fg=teal}
\setbeamercolor{alerted text}{fg=rust}
\setbeamerfont{frametitle}{size=\Large,series=\bfseries}
\setbeamertemplate{navigation symbols}{}
\setbeamertemplate{footline}{\hfill\scriptsize Siyan Wang\enspace|\enspace Preliminary reproduction\enspace|\enspace\insertframenumber/\inserttotalframenumber\hspace{5mm}\vspace{3mm}}
\setbeamersize{text margin left=7mm,text margin right=7mm}
\newcommand{\src}[1]{\par\vspace{1mm}{\fontsize{6.5}{8}\selectfont\color{gray}#1\par}}
\title{9/24 research meeting}
\author{Siyan Wang}
\date{23 September 2026}
\begin{document}

\begin{frame}[plain,noframenumbering]
  \titlepage
\end{frame}


\section{European heatwave attribution: reproduction results}

\begin{frame}[plain,noframenumbering]
    \sectionpage
\end{frame}
% ==============================================================================
% Slide 01 — Data sources and dataset characteristics
% Purpose: classify the four products, retain their temporal coverage, and
% summarize the main strength and limitation of each dataset.
% ==============================================================================
\begin{frame}{Four temperature products: type, coverage and characteristics}
\footnotesize
\renewcommand{\arraystretch}{1.12}
\begin{tabular}{p{1.25cm}p{3.35cm}p{2.35cm}p{6.15cm}}\toprule
Product & Type of data & Temporal coverage & Main characteristics\\\midrule
ERA5 & Reanalysis: observations assimilated into a numerical weather model & 1950--2025 & Continuous and physically consistent; model-dependent.\\
E-OBS & Gridded station observations over European land (spatial interpolation) & 1920--2025 & High European resolution with uncertainty estimates; coverage varies by country and period.\\
CPC & Gridded station-based daily temperature analysis (GTS observations) & 1979--2025 & Frequently updated, but has the shortest record.\\
Berkeley & Observation-based gridded daily land temperature reconstruction & 1897--2023 annual; 1897--2024 June & Longest record with breakpoint correction; daily data are experimental.\\\bottomrule
\end{tabular}
\medskip
\scriptsize All four products are gridded. ERA5 is constrained by a numerical model and data assimilation; the other three are constructed primarily from station observations. Agreement across products therefore tests robustness to different data and processing choices.
\src{Sources: report Appendix A1.1 (p.23); NOAA CPC, E-OBS, ERA5 and Berkeley Earth product documentation; saved fit-year records. Temporal coverage refers to the periods used in our reproduction.}
\end{frame}



% ==============================================================================
% Slide 02 — GMST preprocessing
% Purpose: distinguish the report's smoothing convention from our provisional
% endpoint extrapolation. The GEV model is defined on Slide 03.
% ==============================================================================
\begin{frame}{GMST preprocessing and the 2026 endpoint}
\begin{columns}[T,onlytextwidth]
\column{.51\textwidth}\small
\textbf{Four-year smoothed GMST}\par\smallskip
\footnotesize NASA GISTEMP global land--ocean anomalies (1951--1980 baseline):
\[
g_y=\frac1{12}\sum_{m=1}^{12}g_{y,m},\qquad
G_t=\frac14\sum_{k=-2}^{1}g_{t+k}.
\]
The four-year mean is assigned to year 3 of the window. Thus
\[
G_{2026}=\operatorname{mean}(g_{2024},g_{2025},g_{2026},g_{2027}).
\]
\column{.45\textwidth}\small
\textbf{Our provisional endpoint}\par\smallskip
\footnotesize
\begin{tabular}{rrrr}\toprule
2024 & 2025 & 2026* & 2027*\\\midrule
1.28583 & 1.19333 & 1.22714 & 1.26606\\\bottomrule
\end{tabular}
\par\medskip
\[
\widehat G_{2026}=1.24309^\circ\mathrm C.
\]
* 2026--2027 are extrapolated by OLS over 2011--2025. The report specifies the window, but not how its endpoint values were obtained.
\end{columns}
\medskip\small
The counterfactual world they chose is 1976  (50 years ago, and the year of one of the
most significant heatwaves in the UK in terms of persistently high temperatures), and 2003, when the
first major European heatwave of the century occurred. The climate is $0.6$ and $1.1^\circ$C cooler than now.

\src{Sources: report Appendix A1.1 and Section 1.3; run\_cpc.R; saved GMST endpoint assumptions. Values are anomalies in $^\circ$C and are rounded.}
\end{frame}


% ==============================================================================
% Slide 03 — Statistical model and attribution quantities
% Purpose: define event construction, the GEV model, RR, and intensity change.
% ==============================================================================
\begin{frame}{A location-shift GEV links extremes to GMST}
\begin{columns}[T,onlytextwidth]
\column{.49\textwidth}
\textbf{Event construction}\par\medskip
\small $40$--$60^\circ$N, $10.5^\circ$W--$20^\circ$E.\\[2mm]
Cosine-latitude weighted regional mean; then three-day means and June / annual maxima.\\[2mm]
Historical fitting excludes 2026; complete periods require at least 95\% valid daily values.
\medskip
\begin{tabular}{lrr}\toprule
Current return period & Tx3x & Tn3x\\\midrule
Annual & 25 yr & 20 yr\\
June & 100 yr & 100 yr\\\bottomrule
\end{tabular}
\medskip\par
\scriptsize\textbf{Model choice in the report.} Appendix A1.2 also tests an ENSO/NAO covariate. Table 3 finds no AIC improvement over GMST alone, so these indices are ``not considered for the remainder of this study.'' Our retained results therefore use GMST only.
\column{.47\textwidth}
\textbf{Nonstationary GEV}\par
\footnotesize Each annual extreme $X_t$ is paired with its smoothed GMST $G_t$; one model is fitted across all years.
\[
X_t\sim\mathrm{GEV}(\mu_t,\sigma,\xi),\quad
\mu_t=\mu_0+\alpha G_t.
\]
\scriptsize $\alpha$: location change per $1^\circ$C GMST; $\sigma,\xi$: constant scale and shape. The candidate extension is $\mu_t=\mu_0+\alpha G_t+\beta I_t$, where $I_t$ is detrended Ni\~no3.4 or NAO and $\beta$ is its location effect.\par\smallskip
\textbf{Attribution quantities}\par
$G_F$: 2026 GMST; $G_C=G_F-\Delta G$, where $\Delta G=0.6$ or $1.1^\circ$C. $F_F,F_C$: fitted GEV CDFs; $T$: return period.
\[
x_F=F_F^{-1}(1-1/T),\qquad
RR=\frac{1/T}{1-F_C(x_F)}.
\]
\[
\Delta I=F_F^{-1}(1-1/T)-F_C^{-1}(1-1/T)
=\alpha\Delta G.
\]
\end{columns}
\src{Source: report Appendix A1.2 and Section 2.1.1 (Table 3); saved run scripts and gev.R.}
\end{frame}


% ==============================================================================
% Slide 04 — Figure 7 comparison for ERA5
% Purpose: compare the report panel directly with our ERA5 reproduction.
% ==============================================================================
\begin{frame}{Figure 7 comparison: ERA5 annual Tx3x}
\begin{columns}[T,onlytextwidth]
\column{.47\textwidth}\centering\small\textbf{Report: Figure 7a}\par
\includegraphics[height=4.3cm,trim=104 575 303 72,clip]{0923_reproduction_assets/report_page12.pdf}
\column{.49\textwidth}\centering\small\textbf{Our reproduction: ERA5}\par
\includegraphics[width=\linewidth,height=4.3cm,keepaspectratio,trim=0 263 341 25,clip]{0923_reproduction_assets/figure7.pdf}
\end{columns}
\medskip\small
Both fits show a strong upward shift from the 1976-like climate to the current climate. Our fitted current 25-year level is $30.90^\circ$C; its past return period is about $1.01\times10^6$ years.
\src{Report p.12, Fig.7a; saved ERA5 Figure 7 reconstruction. Our input is daily ERA5 Tmax from KNMI Climate Explorer; the fit uses 1950--2025 and 300 parameter-bootstrap replicates for the plotted bands.}
\end{frame}


% ==============================================================================
% Slide 05 — Figure 7 comparison for CPC
% Purpose: compare the report panel directly with our CPC reproduction.
% ==============================================================================
\begin{frame}{Figure 7 comparison: CPC annual Tx3x}
\begin{columns}[T,onlytextwidth]
\column{.47\textwidth}\centering\small\textbf{Report: Figure 7b}\par
\includegraphics[height=4.3cm,trim=298 575 107 72,clip]{0923_reproduction_assets/report_page12.pdf}
\column{.49\textwidth}\centering\small\textbf{Our reproduction: CPC}\par
\includegraphics[width=\linewidth,height=4.3cm,keepaspectratio,trim=346 263 0 25,clip]{0923_reproduction_assets/figure7.pdf}
\end{columns}
\medskip\small
Both show an upward shift from 1976-like (blue) to current climate (red). Our fitted current 25-year level is $33.18^\circ$C; its past return period is 229 years.

\src{Report p.12, Fig.7b (original crop); saved Figure 7 reconstruction. Purple line: current 25-year threshold. Our bands: 300 parameter-bootstrap replicates; adjusted observations shown as points.}
\end{frame}



% ==============================================================================
% Slide 06 — Figure 7 comparison for Berkeley
% Purpose: compare the report panel directly with our Berkeley reproduction.
% ==============================================================================
\begin{frame}{Figure 7 comparison: Berkeley annual Tx3x}
\begin{columns}[T,onlytextwidth]
\column{.47\textwidth}\centering\small\textbf{Report}\par
\includegraphics[height=4.3cm,trim=104 390 291 264,clip]{0923_reproduction_assets/report_page12.pdf}
\column{.49\textwidth}\centering\small\textbf{Our reproduction}\par
\includegraphics[width=\linewidth,height=4.3cm,keepaspectratio,trim=0 0 341 277,clip]{0923_reproduction_assets/figure7.pdf}
\end{columns}
\medskip\small Our current 25-year threshold is $31.92^\circ$C; the fitted 1976-like return period is 4,876 years. Both plots show warming, but their threshold levels and tails differ.
\medskip\par\footnotesize 
\src{Sources: report p.12, Figure 7; saved Figure 7 reconstruction and summary. Original axes retained; panel scales differ. Our confidence bands use 300 parameter-bootstrap replicates.}
\end{frame}



% ==============================================================================
% Slide 07 — Figure 7 comparison for E-OBS
% Purpose: compare the report panel directly with our E-OBS reproduction.
% ==============================================================================
\begin{frame}{Figure 7 comparison: E-OBS annual Tx3x}
\begin{columns}[T,onlytextwidth]
\column{.47\textwidth}\centering\small\textbf{Report}\par
\includegraphics[height=4.3cm,trim=298 390 107 264,clip]{0923_reproduction_assets/report_page12.pdf}
\column{.49\textwidth}\centering\small\textbf{Our reproduction}\par
\includegraphics[width=\linewidth,height=4.3cm,keepaspectratio,trim=346 0 0 277,clip]{0923_reproduction_assets/figure7.pdf}
\end{columns}
\medskip\small Our current 25-year threshold is $32.05^\circ$C; the fitted 1976-like return period is 52,044 years. Both plots show warming, but their threshold levels and tails differ.
\medskip\par\footnotesize The report uses v32.0e; our input uses v33e.
\src{Sources: report p.12, Figure 7; saved Figure 7 reconstruction and summary. Original axes retained; panel scales differ. Our confidence bands use 300 parameter-bootstrap replicates.}
\end{frame}



% ==============================================================================
% Slide 08 — Original Table 2 benchmark
% Purpose: show the synthesized numerical target before our comparisons.
% ==============================================================================
\begin{frame}{Table 2: the report's numerical target}
\centering
\includegraphics[width=13.4cm,height=4.7cm,keepaspectratio,trim=44 92 59 528,clip]{0923_reproduction_assets/report_page12.pdf}
\par\medskip\raggedright
\small The report synthesizes results across observational and reanalysis datasets. Our current outputs are \textbf{individual-dataset estimates}; the synthesis is still pending.
\medskip\par\footnotesize Parentheses contain reported 95\% intervals. ``$>1$m'' means $>10^6$. Annual return periods: 25 yr for Tx3x, 20 yr for Tn3x; June: 100 yr for both.
\src{Original Table 2, report p.12; caption and synthesis description on p.13. The next slides use this table as a benchmark, not as an individual-dataset ground truth.}
\end{frame}



% ==============================================================================
% Slide 09 — Table 2 intensity comparison
% Purpose: compare synthesis with individual-dataset intensity estimates.
% ==============================================================================
\begin{frame}{Table 2 comparison: change in intensity since 2003}
\centering\includegraphics[width=14.2cm,height=5.25cm,keepaspectratio]{0923_reproduction_assets/intensity_comparison.pdf}
\par\raggedright\small
Our estimates share the report's positive warming direction. Dataset-level differences from the synthesis remain substantial; interval overlap is not a test of reproduction accuracy.
\src{Black: report Table 2 synthesis; colors: saved single-dataset estimates. Bars: 95\% intervals. Our intervals use 200 requested parameter-bootstrap fits (173--200 successful). Report values transcribed directly. Units: $^\circ$C.}
\end{frame}



% ==============================================================================
% Slide 10 — Table 2 probability-ratio comparison
% Purpose: compare synthesis with individual-dataset PR/RR estimates.
% ==============================================================================
\begin{frame}{Table 2 comparison: probability ratios}
\small Each cell shows \textbf{2003-like / 1976-like}. PR in the report and RR in our code denote the same probability ratio.
\medskip
\scriptsize
\setlength{\tabcolsep}{3pt}
\renewcommand{\arraystretch}{1.5}
\centering
\begin{tabular}{llrrrrr}\toprule
Index & Period & Report & ERA5 & CPC & E-OBS & Berkeley\\\midrule
Tx3x & Annual & $13.9/580$ & $37.1/40221$ & $3.74/9.15$ & $19.4/2082$ & $9.53/195$\\
Tx3x & June & $286/{>}10^6$ & $3788/\infty$ & $\infty/\infty$ & $77.3/3.09{\times}10^{17}$ & $57.3/\infty$\\
Tn3x & Annual & $163/{>}10^6$ & --- & $103/6.47{\times}10^{10}$ & $7.68/255$ & ---\\
Tn3x & June & $295/{>}10^6$ & --- & $\infty/\infty$ & $7.72/61.6$ & ---\\\bottomrule
\end{tabular}
\par\medskip\raggedright\footnotesize
Point estimates only. For Annual Tx3x vs 2003-like, 95\% intervals are: report [2.28, 4980], ERA5 [5.61, $\infty$], CPC [1.63, 56.0], E-OBS [4.51, $\infty$], Berkeley [3.52, 15505].\par\medskip
\small Large discrepancies in tail probabilities remain. $\infty$ indicates zero fitted counterfactual probability at the threshold; the report's $>10^6$ is a reporting bound, not an exact infinity.
\src{Sources: original Table 2 and saved PROVISIONAL result CSVs. Same prescribed return periods and GMST decrements; our input snapshots, preprocessing and synthesis are not fully aligned.}
\end{frame}



% ==============================================================================
% Slide 11 — Reproduction status and remaining gaps
% Purpose: summarize agreements, discrepancies, and unfinished work.
% ==============================================================================
\begin{frame}{Reproduction assessment and remaining gaps}
\begin{columns}[T,onlytextwidth]
\column{.48\textwidth}
\textbf{What we reproduced}\par\medskip
\small The location-shift GEV workflow, fixed return periods and $0.6/1.1^\circ$C GMST comparisons.\\[3mm]
ERA5 Tx3x, CPC, E-OBS and Berkeley individual fits, plus the structure of Figure 7.\\[3mm]
The report's qualitative warming direction. Quantitative agreement is incomplete, especially for risk ratios.
\column{.48\textwidth}
\textbf{What prevents an exact match}\par\medskip
\small E-OBS version and source snapshots differ; land masks, missing-data rules and time windows need alignment.\\[2mm]
Our 2026 GMST endpoint is extrapolated; the author's endpoint inputs have not been matched.\\[2mm]
Bootstrap design and optimizer may differ. Our intervals omit GMST uncertainty.
\end{columns}
\medskip\small\textcolor{teal}{ Next Step: Checking/considering if the  ENSO/NAO covariate are not improving the model}
\src{Evidence: report Appendix A1; original Figures / Table 2; saved reproduction scripts and output CSVs. The present comparisons do not constitute a completed numerical replication.}
\end{frame}

% ==============================================================================
% Section 2 PN,PS,PNS in Hannart et al. (2016)
% ==============================================================================

\section{PN,PS,PNS in Hannart et al. (2016)}

\begin{frame}[plain,noframenumbering]
    \sectionpage
\end{frame}


% Slide 12: Establish the factual/counterfactual notation before introducing
% conditional probabilities. Y_0 and Y_1 refer to the same situation.
\begin{frame}{Set-up: two potential outcomes}
  \begin{itemize}
    \item $X=1$: the forcing is present; $X=0$: it is absent.
    \item $Y=1$: the defined extreme event occurs.
    \item $Y_1$ and $Y_0$: whether the event would occur with and without the forcing, respectively, for the same situation.
  \end{itemize}
  \begin{block}{Two questions}
    For an event that occurred, would it have occurred without the forcing?\\[0.4em]
    For an event that did not occur, would adding the forcing make it occur?
  \end{block}
  \src{Source: Hannart et al. (2016).}
\end{frame}

% Slide 13: These are definitions, not formulas inferred from p_0 and p_1.
% Consistency gives Y=Y_1 for X=1 and Y=Y_0 for X=0.
% PN conditions on an observed event; PS conditions on an observed non-event.
% PNS is unconditional over the target population.
\begin{frame}{Three probabilities of causation}
  \begin{align*}
    \mathrm{PN}  &= P(Y_0=0\mid X=1,Y=1),\\[0.35em]
    \mathrm{PS}  &= P(Y_1=1\mid X=0,Y=0),\\[0.35em]
    \mathrm{PNS} &= P(Y_0=0,Y_1=1).
  \end{align*}
  \begin{itemize}
    \item \textbf{PN (necessity):} Among events that occurred with forcing, how many would disappear without it?
    \item \textbf{PS (sufficiency):} Among non-events without forcing, how many would occur if forcing were added?
    \item \textbf{PNS (necessity and sufficiency):} How many situations are switched from $0$ to $1$ by the forcing?
  \end{itemize}
  \src{Source: Hannart et al. (2016).}
\end{frame}

% Slide 14: Enumerate all four possible pairs of potential outcomes.
% Monotonicity is an individual-level statement: it excludes (Y_0,Y_1)=(1,0).
% A population-level finding p_1 >= p_0 alone does not establish monotonicity.
% Exogeneity allows observed risks P(Y=1|X=x) to represent P(Y_x=1).
\begin{frame}{Why the formulas work: response types}
  \centering
  \begin{tabular}{ccc}
    \toprule
    $Y_0$ & $Y_1$ & Response to the forcing \\
    \midrule
    0 & 0 & Never occurs \\
    0 & 1 & Occurs only with forcing (caused) \\
    1 & 0 & Prevented by forcing \\
    1 & 1 & Occurs either way \\
    \bottomrule
  \end{tabular}
  \vspace{0.8em}

  \begin{block}{Assumptions for exact identification}
    \textbf{Monotonicity:} $Y_1\geq Y_0$ for every situation, excluding the $(1,0)$ row.\\
    \textbf{Exogeneity:} the forcing is independent of potential outcomes, so the factual and counterfactual risks equal $p_x=P(Y_x=1)$.
  \end{block}
  \src{Source: Hannart et al. (2016).}
\end{frame}

% Slide 15: Under monotonicity, p_0 is the share of (1,1) cases, while p_1
% includes both (1,1) and (0,1). Their difference is the share of (0,1).
% With exogeneity, conditioning on X does not change the share of response
% types, so PN and PS use the same (0,1) numerator with different denominators.
% The denominators must be positive for the corresponding conditional
% probabilities to be defined.
\begin{frame}{Derivation from the two risks}
  Let $p_1=P(Y_1=1)$ and $p_0=P(Y_0=1)$. Under monotonicity,
  \[
    \underbrace{p_1}_{(0,1)+(1,1)}
    -\underbrace{p_0}_{(1,1)}
    =\underbrace{P(Y_0=0,Y_1=1)}_{\text{caused type}}.
  \]
  Under exogeneity as well, the caused type is the numerator of each conditional probability:
  \begin{align*}
    \mathrm{PNS} &= p_1-p_0,\\
    \mathrm{PN}  &= \frac{p_1-p_0}{p_1}
                  =1-\frac{p_0}{p_1}\quad(p_1>0),\\
    \mathrm{PS}  &= \frac{p_1-p_0}{1-p_0}
                  =1-\frac{1-p_1}{1-p_0}\quad(p_0<1).
  \end{align*}
  \src{Source: Hannart et al. (2016).}
\end{frame}

% Slide 16: Connect PN to the familiar fraction of attributable risk (FAR).
% This equality needs both exogeneity and monotonicity. Without monotonicity,
% p_1-p_0 is the caused share minus the prevented share; the displayed
% expressions are then lower bounds rather than exact probabilities.
% The example highlights that PN and PS can differ because they condition
% on different observed groups.
\begin{frame}{Interpretation and limits}
  \begin{itemize}
    \item \textbf{Ex post attribution:} PN asks whether an observed event needed the forcing. Under the two assumptions, $\mathrm{PN}=1-p_0/p_1=\mathrm{FAR}$.
    \item \textbf{Prospective question:} PS asks whether adding the forcing would produce an event that otherwise would not occur.
    \item \textbf{Without monotonicity:} $p_1-p_0$ subtracts the fraction prevented by forcing from the fraction caused by it. The probabilities are generally not identified by $p_0,p_1$ alone; the expressions above give lower bounds (truncated at zero).
  \end{itemize}
  \begin{block}{Example: $p_0=0.02$, $p_1=0.10$}
    $\mathrm{PNS}=0.08,\qquad \mathrm{PN}=0.80,\qquad \mathrm{PS}\approx0.082.$
  \end{block}
  \src{Source: Hannart et al. (2016).}
\end{frame}

\end{document}


